% META: {"id": "TSPE-SUITLIM-ME-001", "chapitre": "TSPE-SUITES-LIMITES", "type_objet": "methode", "capacites_codes": ["C1"], "methodes": ["M1"], "sources_inspiration": [], "status": "approved"}
\begin{fichemethode}{M1}{Déterminer la limite d'une suite}

\quandUtiliser{%
  Utiliser cette méthode lorsque l'énoncé demande de « calculer la limite », « étudier la convergence » ou « déterminer le comportement en $+\infty$ » d'une suite.
}

\pasApas{%
  \begin{enumerate}
    \item \textbf{Identifier la forme de la suite :}
      \begin{itemize}
        \item Suite explicite $u_n = f(n)$ : utiliser les règles d'opérations sur les limites.
        \item Suite géométrique $u_n = a \cdot q^n$ : utiliser le tableau des limites de $q^n$.
        \item Suite définie par récurrence $u_{n+1} = f(u_n)$ : étudier d'abord la convergence, puis chercher le point fixe.
      \end{itemize}
    \item \textbf{Vérifier s'il y a une forme indeterminee} (type $\dfrac{\infty}{\infty}$, $+\infty - \infty$, $0 \times \infty$).
    \item \textbf{Si forme indeterminee :} factoriser par le terme dominant, utiliser un encadrement, ou une suite auxiliaire.
    \item \textbf{Conclure :} enoncer la limite et preciser si la suite converge ou diverge.
  \end{enumerate}
}

\exempleRedige{%
  \textbf{Exemple.} Déterminer la limite de $u_n = \dfrac{3n^2 + n}{n^2 - 1}$.

  \medskip
  \commentaireMarge{Forme $\frac{\infty}{\infty}$ : factoriser par $n^2$.}%
  On factorise par $n^2$ au numérateur et au dénominateur :
  \[
    u_n = \frac{n^2(3 + 1/n)}{n^2(1 - 1/n^2)} = \frac{3 + 1/n}{1 - 1/n^2}.
  \]

  \commentaireMarge{Appliquer les opérations sur les limites.}%
  Or $\lim \dfrac{1}{n} = 0$ et $\lim \dfrac{1}{n^2} = 0$, donc :
  \[
    \lim_{n \to +\infty} u_n = \frac{3 + 0}{1 - 0} = 3.
  \]

  La suite $(u_n)$ converge vers $3$.
}

\pieges{%
  \begin{itemize}
    \item \textbf{Piège 1.} Écrire directement $\dfrac{\infty}{\infty} = 1$. C'est une forme indeterminee ; il faut factoriser.
    \item \textbf{Piège 2.} Pour une suite recurrente, calculer le point fixe sans avoir prouve la convergence. Le point fixe ne donne la limite que \emph{si} la suite converge.
  \end{itemize}
}

\verifier{%
  Comparer la limite trouvee avec les premieres valeurs numériques. Si $u_n \to 3$ mais $u_{10} = 500$, il y a une erreur.
}

\sEntrainer{\refExos{M1}}

\end{fichemethode}
